\begin{statementbox}
固定中心 $B(a,b)$，点 $P(x_1,y_1)$ 绕 $B$ 旋转有向角 $\theta$ 到 $Q(x_2,y_2)$。
\textbf{逆时针取正，顺时针取负；旋转保持 $BP=BQ$。}
令 $u=x_1-a,\ v=y_1-b$，即 $\overrightarrow{BP}$ 的相对坐标，则
\[
\boxed{\begin{cases}
x_2=a+u\cos\theta-v\sin\theta,\\
y_2=b+u\sin\theta+v\cos\theta.
\end{cases}}
\]
\begin{center}\begin{tikzpicture}[x=1cm,y=1cm,>=Stealth,font=\small]
\path[use as bounding box] (-6.6,-0.6) rectangle (6.6,4.9);

\begin{scope}[shift={(-2.7,1.1)},scale=0.78]
\draw[->,gray] (-1,0)--(6.4,0) node[right] {$x$};
\draw[->,gray] (0,-0.5)--(0,4.5) node[above] {$y$};
\coordinate (B) at (2,1); \coordinate (P) at (5,2);
\coordinate (Q) at ({2+sqrt(2)},{1+2*sqrt(2)});
\draw[blue!75!black,very thick,->] (B)--(P);
\draw[teal!75!black,very thick,->] (B)--(Q);
\draw[gray,dashed] (B)--(5,1)--(P);
\node[below] at (3.5,1) {$u=x_1-a$};
\node[right] at (5,1.5) {$v$};
\draw[orange!85!black,thick,->] ($(B)+(18.435:1.05)$) arc (18.435:63.435:1.05);
\node[orange!85!black] at (3.45,2.35) {$\theta$};
\fill (B) circle(2pt) node[below left] {$B(a,b)$};
\fill[blue!75!black] (P) circle(2pt) node[above right] {$P(x_1,y_1)$};
\fill[teal!75!black] (Q) circle(2pt) node[above] {$Q(x_2,y_2)$};
\node[below left] at (0,0) {$O$};
\end{scope}
\node[align=left,anchor=east] at (-3.05,3.15) {旋转中心是 $B$\\不是原点 $O$};
\node[align=left,anchor=west] at (3.6,3.7) {$BP=BQ$\\逆时针：$\theta>0$\\顺时针：$\theta<0$};
\node at (0,0) {先看 $\overrightarrow{BP}=(u,v)$，旋转后再加回 $(a,b)$};

\end{tikzpicture}\end{center}

\textbf{读图：} 转动的是从 $B$ 出发的箭头，中心 $B$ 始终不动。图中取
$\overrightarrow{BP}=(3,1),\ \theta=45^\circ$；一般公式适用于任意实数 $u,v$ 和任意角。

\textbf{两个边界：} $B=O$ 时令 $a=b=0$；$P=B$ 时旋转后仍有 $Q=B$，此时不能由对应点反求角度。
\end{statementbox}
\clearpage
\begin{statementbox}
\textbf{顺时针也可用正角 $\beta$ 表示。} 若顺时针旋转 $\beta>0$，代 $\theta=-\beta$，得
\[
\boxed{\begin{cases}
x_2=a+u\cos\beta+v\sin\beta,\\
y_2=b-u\sin\beta+v\cos\beta.
\end{cases}}
\]
\textbf{特殊角速查（表内均为相对坐标）：}
\[
\begin{array}{c|c}
\theta&(u,v)\text{ 变为 }(s,t)\\\hline
45^\circ&((u-v)/\sqrt2,(u+v)/\sqrt2)\\[3pt]
90^\circ&(-v,u)\\[3pt]
-90^\circ&(v,-u)\\[3pt]
180^\circ&(-u,-v)
\end{array}
\]
其中 $s=x_2-a,t=y_2-b$。$90^\circ$ 是 V039 的垂直等长特例；$180^\circ$ 表示 $B$ 为 $PQ$ 的中点。

\textbf{反求角度的条件与结论：} 先确认 $BP=BQ=r>0$，再用
\[
\boxed{\cos\theta=\frac{us+vt}{r^2},\qquad
\sin\theta=\frac{ut-vs}{r^2}.}
\]
这两式只能确定相差整周的角；若需唯一角度，应指定角度范围。例如取
$-180^\circ<\theta\le180^\circ$ 时，正角表示逆时针、负角表示顺时针的最短转动；$180^\circ$ 两种方向到达同一点，$0^\circ$ 表示点未改变。
\end{statementbox}
